\begin{explanationbox}
\textbf{\textcolor{CExplanation}{一句话直觉}}
\textbf{圆的直径所对圆周角是直角，圆变瘦成椭圆后，两条连线斜率被同倍压缩，于是斜率的乘积仍为固定负数。}

\medskip
\begin{minipage}{\linewidth}
\textbf{\textcolor{CExplanation}{最低门槛：先想圆，再想压缩（逐步理解坐标怎么变）}}
\begin{center}
\begin{tikzpicture}[font=\small]
 \path[use as bounding box] (-2.8,-3.9) rectangle (9.8,3.3);
 \foreach \shift/\ry/\name in {0/2/圆,7/1.2/椭圆}{
 \begin{scope}[shift={(\shift,0)}]
 \node at (0,2.95) {\name：横坐标保持不变};
 \draw[->,gray] (-2.4,0)--(2.4,0);
 \draw[blue!60!black,thick] (0,0) ellipse (2 and \ry);
 \coordinate (L) at (-2,0);\coordinate (R) at (2,0);\coordinate (T) at (1.2,{.8*\ry});
 \draw[teal!70!black,thick] (L)--(T);
 \draw[magenta!70!black,thick] (R)--(T);
 \draw[dashed] (1.2,0)--(T);
 \fill (T) circle (1.5pt);
 \node[below left] at (L) {$A_1$};\node[below right] at (R) {$A_2$};
 \end{scope}}
 % Explicit right angle on the circle, with its own coordinates.
 \coordinate (CL) at (-2,0);\coordinate (CR) at (2,0);\coordinate (Q) at (1.2,1.6);
 \pic[draw,angle radius=3mm] {right angle=CL--Q--CR};
 \node[above right] at (Q) {$Q(u,v)$};
 \node[above right] at (8.2,.96) {$P(u,\frac ba v)$};
 \node at (0,-2.6) {$k_{QA_1}k_{QA_2}=-1$};
 \node at (7,-2.6) {$k_{PA_1}k_{PA_2}=-\bigl(\frac ba\bigr)^2$};
 \draw[->,thick] (2.65,1.95)--(4.35,1.95) node[midway,above] {$y\mapsto\frac ba y$};
 \node at (3.5,-3.4) {图取 $a=2,b=1.2$，两侧单位长度相同};
\end{tikzpicture}
\end{center}

\end{minipage}
\textbf{先认清选图的理由：}选圆 $C_0:u^2+v^2=a^2$，半径恰好等于椭圆的长半轴 $a$；
圆与椭圆的左右端点都是 $A_1(-a,0),A_2(a,0)$，因此可以让这两个端点固定不动。
图中两侧\textbf{单位长度一致}；此处采用的映射是\textbf{只改高度、不改横坐标}，而不是随便把两条曲线上的点配对。

\begin{enumerate}[leftmargin=*,itemsep=0.45em,topsep=0.35em]
 \item \textbf{圆上选一个点。}设 $Q=(u,v)$，则 $u^2+v^2=a^2$。
  因为圆上的点只有两种高度 $v=\pm\sqrt{a^2-u^2}$，先记住它位于横轴上方还是下方。
 \item \textbf{先固定横坐标，再问目标点有多高。}让对应的椭圆点写成 $P=(u,w)$。
  既然 $P$ 要落在 $\frac{x^2}{a^2}+\frac{y^2}{b^2}=1$ 上，就必须有
  \[
   \frac{u^2}{a^2}+\frac{w^2}{b^2}=1
   \quad\Longrightarrow\quad
   w^2=\frac{b^2}{a^2}(a^2-u^2)=\frac{b^2}{a^2}v^2.
  \]
 \item \textbf{为什么选正的比例？}我们要求上下位置不翻转：$Q$ 在上半圆，$P$ 就在上半椭圆；
  $Q$ 在下半圆，$P$ 也在下半椭圆。于是 $w$ 与 $v$ 同号：
  \[
    \boxed{w=\frac ba v},\qquad \boxed{Q(u,v)\mapsto P\left(u,\frac ba v\right)}.
  \]
  这说明纵坐标乘 $b/a$ \textbf{不是凭记忆硬指定}；它是由\textbf{同一横坐标、落在椭圆上、上下位置不翻转}确定的。
 \item \textbf{再检查整个图形。}圆的最高点 $(0,a)$ 变成椭圆最高点 $(0,b)$；
  最低点 $(0,-a)$ 变成 $(0,-b)$；而横轴上 $v=0$ 的点完全不动。
  因 $0<b/a<1$，所有高度等比缩小，圆被纵向压成椭圆。
 \item \textbf{反方向：椭圆拉伸回圆。}若先有 $P=(x,y)$ 在椭圆上，
  令 $Q=(x,\frac ab y)$，则
  \[
    x^2+\left(\frac ab y\right)^2=a^2.
  \]
  所以 $Q$ 必在半径 $a$ 的圆上。\textbf{压缩与拉伸互为逆操作}，不是只能由圆推椭圆。
\end{enumerate}

\medskip
\textbf{\textcolor{CExplanation}{圆与椭圆的双向坐标对应}}
下面两图将圆与椭圆叠放在\textbf{同一坐标系}中：灰色虚线是圆，蓝色实线是椭圆。
橙色箭头表示对应点沿同一条竖直线移动；上下两图的坐标单位、曲线和取点完全相同。
% Equal units on both axes; a=3, b=1.8, Q=(1.8,2.4), P=(1.8,1.44).
\begin{center}
\begin{tikzpicture}[font=\small,>=Stealth]
 \path[use as bounding box] (-4,-4.6) rectangle (6.1,4.2);
 \node at (.8,3.85) {圆与椭圆的双向坐标对应（一）：纵向压缩};
 \draw[gray,thick,dashed] (0,0) circle (3);
 \draw[blue!70!black,very thick] (0,0) ellipse (3 and 1.8);
 \draw[->,gray] (-3.6,0)--(3.6,0) node[right] {$x$};
 \coordinate (A) at (-3,0);\coordinate (B) at (3,0);
 \coordinate (Q) at (1.8,2.4);\coordinate (P) at (1.8,1.44);\coordinate (H) at (1.8,0);
 \draw[gray!60,thick] (A)--(Q)--(B);
 \draw[teal!70!black,very thick] (A)--(P);
 \draw[magenta!70!black,very thick] (P)--(B);
 \draw[dashed,gray,thick] (H)--(P);
 \draw[->,orange!90!black,line width=1.6pt] (Q)--(P);
 \draw[gray,thin] (Q)--(3.7,2.4);
 \draw[magenta!70!black,thin] (P)--(3.7,1.44);
 \node[right] at (3.8,2.4) {$Q(u,v)$};
 \node[right] at (3.8,1.44) {$P(u,\frac ba v)$};
 \node[right,orange!90!black] at (3.8,.55) {纵坐标乘 $\frac ba<1$};
 \fill[blue!70!black] (A) circle (2pt) node[below left] {$A_1(-a,0)$};
 \fill[blue!70!black] (B) circle (2pt) node[below right] {$A_2(a,0)$};
 \fill[gray] (Q) circle (2.2pt);
 \fill[magenta!70!black] (P) circle (2.2pt);
 \fill[gray] (H) circle (1.7pt) node[below] {$H(u,0)$};
 \node[below left] at (0,0) {$O$};
 \node at (.6,-3.55) {$u^2+v^2=a^2\quad\xrightarrow{\ T(u,v)=(u,\frac ba v)\ }\quad
  \frac{x^2}{a^2}+\frac{y^2}{b^2}=1$};
 \node at (.6,-4.2) {横坐标不变，端点不动；上、下半平面分别对应（$a>b>0$）};
\end{tikzpicture}
\end{center}

% Inverse of C064_circle_to_ellipse; identical curves and units.
\begin{center}
\begin{tikzpicture}[font=\small,>=Stealth]
 \path[use as bounding box] (-4,-4.6) rectangle (6.1,4.2);
 \node at (.8,3.85) {圆与椭圆的双向坐标对应（二）：反向还原};
 \draw[gray,thick,dashed] (0,0) circle (3);
 \draw[blue!70!black,very thick] (0,0) ellipse (3 and 1.8);
 \draw[->,gray] (-3.6,0)--(3.6,0) node[right] {$x$};
 \coordinate (A) at (-3,0);\coordinate (B) at (3,0);
 \coordinate (Q) at (1.8,2.4);\coordinate (P) at (1.8,1.44);\coordinate (H) at (1.8,0);
 \draw[gray!60,thick] (A)--(Q)--(B);
 \draw[teal!70!black,very thick] (A)--(P);
 \draw[magenta!70!black,very thick] (P)--(B);
 \draw[dashed,gray,thick] (H)--(P);
 \draw[->,orange!90!black,line width=1.6pt] (P)--(Q);
 \draw[gray,thin] (Q)--(3.7,2.4);
 \draw[magenta!70!black,thin] (P)--(3.7,1.44);
 \node[right] at (3.8,2.4) {$Q(x,\frac ab y)$};
 \node[right] at (3.8,1.44) {$P(x,y)$};
 \node[right,orange!90!black] at (3.8,.55) {纵坐标乘 $\frac ab>1$};
 \fill[blue!70!black] (A) circle (2pt) node[below left] {$A_1(-a,0)$};
 \fill[blue!70!black] (B) circle (2pt) node[below right] {$A_2(a,0)$};
 \fill[gray] (Q) circle (2.2pt);
 \fill[magenta!70!black] (P) circle (2.2pt);
 \fill[gray] (H) circle (1.7pt) node[below] {$H(x,0)$};
 \node[below left] at (0,0) {$O$};
 \node at (.6,-3.55) {$\frac{x^2}{a^2}+\frac{y^2}{b^2}=1\quad
  \xrightarrow{\ T^{-1}(x,y)=(x,\frac ab y)\ }\quad x^2+\left(\frac ab y\right)^2=a^2$};
 \node at (.6,-4.2) {横坐标不变，端点不动；先压缩再拉伸，回到原来的点};
\end{tikzpicture}
\end{center}

两方向的纵坐标倍率互为倒数：$(b/a)(a/b)=1$。图取 $a=3,b=1.8$，
圆点 $(1.8,2.4)$ 压缩为椭圆点 $(1.8,1.44)$，再拉伸即可还原。

\textbf{变化与不变分别是什么？}令垂足 $H=(u,0)$，则
\[
 A_1,O,A_2,H\ \text{的位置不变},\qquad
 x_P=x_Q=u,\qquad
 HP=\frac ba\,HQ\ \text{（长度）},\qquad
 \frac{y_P}{b}=\frac{v_Q}{a}.
\]
两条从固定端点连向动点的线段\textbf{横向跨度不变}、\textbf{竖向增量乘 $b/a$}；
所以每条斜率都乘 $b/a$，斜率积由 $-1$ 变为 $-(b/a)^2$。
\textbf{特别小心：}这种映射保留直线和共线关系，但通常不保留长度、角度与垂直关系；
圆上的直角标记\textbf{不能}直接搬到椭圆上。


\medskip
\textbf{\textcolor{CExplanation}{为什么是负数？为什么是定值？}}
设 $P(x,y)$，由于 $-a<x<a$，从 $A_1$ 看 $P$ 的横向增量 $x+a>0$，
从 $A_2$ 看则为 $x-a<0$，两个斜率的分子都为 $y$，
因此它们的符号相反。椭圆方程还保证
\[
 \underbrace{y^2}_{\text{两个斜率的分子乘积}}
 =\frac{b^2}{a^2}\underbrace{(a^2-x^2)}_{\text{分母平方差的相反数}}.
\]
正好约掉动点的坐标，只剩负的轴长平方比。

\medskip
\textbf{\textcolor{CExplanation}{高价值识别信号}}
\begin{itemize}[leftmargin=*,itemsep=0.3em,topsep=0.3em]
 \item \textbf{看到“椭圆上一动点 $P$ + 长轴两顶点 + 两条连线斜率”}：直接用斜率积，无须联立两条直线。
 \item \textbf{看到“斜率积 $K$ 反求离心率”}：长轴在 $x$ 轴时 $K=e^2-1$，所以 $e=\sqrt{1+K}$（必有 $-1<K<0$）。
 \item \textbf{看到“以两动点为直径的圆过定点”}：尝试把两动点纵坐标的乘积转成 $k_1k_2$，再令圆方程中可变的一次项消失。
 \item \textbf{看到“固定两点，斜率积是负常数”}：可逆推出椭圆轨迹，并判断固定两点到底是长轴端点还是短轴端点。
\end{itemize}

\medskip
\textbf{\textcolor{CExplanation}{方向统一：看斜率定义，而不是死记字母}}
\begin{itemize}[leftmargin=*,itemsep=0.3em,topsep=0.3em]
 \item 长轴水平时，斜率是 $\frac{y}{x\pm a}$，因此\textbf{横向端点的斜率积}为 $-b^2/a^2$。
 \item 长轴竖直时，斜率是 $\frac{y\pm a}{x}$，因此\textbf{竖向端点的斜率积}为 $-a^2/b^2$。
 \item \textbf{两式互为倒数，不是简单地把“焦点”换成“顶点”。} 注意上下端点的斜率分母是 $x$。
\end{itemize}

\medskip
\textbf{\textcolor{CExplanation}{迁移边界}}
这个结论依赖于\textbf{指定端点和坐标斜率}。平移椭圆（端点也随之平移）不改变斜率积；
若把椭圆任意旋转而仍使用原 $x,y$ 坐标计算斜率，一般不能原样套公式。

\medskip
\textbf{\textcolor{CExplanation}{更深一层：为什么说是仿射缩放，不是相似缩放？}}
\textbf{相似缩放}会让横、纵两个方向乘\textbf{同一个}比例，因此角度保持；
本题的 $T(u,v)=(u,\frac ba v)$ 只改纵向，不改横向，是\textbf{不等比缩放（仿射变换）}。
可以把圆点参数写成 $Q=(a\cos t,a\sin t)$，映射后
$P=(a\cos t,b\sin t)$。两个点保留同一个参数 $t$，但\textbf{$t$ 通常不是点 $P$ 的极角}。
这帮助我们抓住真正的不变量：\textbf{归一化坐标} $x/a=\cos t$、$y/b=\sin t$，而不是图上看到的角度。
\end{explanationbox}
